\section{Conclusion}

In this work, we presented a fully automated, black-box pipeline for detecting unverbalized biases in LLMs: concepts that systematically influence model decisions without being verbalized in chain-of-thought reasoning. Our approach combines LLM-based concept hypothesis generation, counterfactual input variations, and statistical testing with early stopping to efficiently discover biases without requiring predefined categories or manual dataset curation.

On three decision tasks across seven models, our pipeline successfully rediscovered biases that prior work identified manually (gender and race in hiring) while also uncovering new ones not covered by manual analysis (Spanish fluency, English proficiency, writing formality). The cross-task consistency of detected biases, particularly for gender and race, suggests these reflect genuine model behavior rather than task-specific artifacts. The pipeline also correctly filters concepts that are explicitly verbalized in model reasoning, providing evidence for the absence of certain biases. Models also differ in how freely they mention demographic factors. Grok 4.1 Fast acknowledges demographic characteristics in its chain-of-thought far more often than other models, resulting in fewer detected unverbalized biases despite having similar underlying tendencies.

Future work could improve the pipeline by refining verbalization detection with adaptive thresholds or semantic matching, and incorporating evolutionary algorithms to iteratively explore the concept space. Combining automated hypothesis generation with domain expert input may further improve coverage of subtle biases.

As LLMs are increasingly deployed in high-stakes decision-making, understanding the factors that influence their outputs beyond stated reasoning becomes critical. We hope this work contributes to more robust monitoring of model behavior and encourages further investigation into the faithfulness of LLM reasoning.